% language=us runpath=texruns:manuals/ontarget

\startcomponent ontarget-rules

\environment ontarget-style

\startchapter[title={Getting rid of math rules}]

There are a few characteristics of typesetting math that determine how the
internals of the \TEX\ math engine are shaped:

\startitemize
    \startitem
        A formula is made from atoms: in $x + 1$ there are three of them.
        Sometimes a few are combined into subformulas.
    \stopitem
    \startitem
        An atom can have scripts attached: in $x^2_i$ we have a super- and
        subscript. We can also have prescripts and indices.
    \stopitem
    \startitem
        Spacing between atoms depends on what it represents: a operator, binary,
        ordinal etc. Sometimes we look back, sometimes we look ahead.
    \stopitem
    \startitem
        The raw math list is stepwise converted to a normal node list. A next pass
        can use information from a previous one.
    \stopitem
    \startitem
        Fractions are implemented as dedicated objects with two atoms on top of
        each other and optionally scripts and fences.
    \stopitem
    \startitem
        Accents can be put on top of or below an atom or subformula and are an
        object too. Again they can have scripts.
    \stopitem
    \startitem
        Fences that are supposed to scale can wrap an atom or subformula and are
        an object. Scripts can be attached as usual.
    \stopitem
    \startitem
        Radicals have a symbols in from that extends over the atom or subformula
        and therefore they are objects. Such an object can have a degree and
        scripts.
    \stopitem
\stopitemize

Some specific constructs like integrals and actuarians use these mechanisms and
don't have their own. So, for instance \quote {radical} is more a generic
construct than one specific for roots.

Although in principle one can implement a math renderer in \TEX\ it is the
interplay between atoms, objects, scripts, fence sizing, spacing, etc,\ that
makes the argument for using a hard|-|coded solution. Because there is some
general agreement about how math should be rendered, this also results in some
standardization. There are plenty parameters that drive this process and
especially the extensive interaction between fonts and \quotation {what to do in
what situation} makes it more convenient to do this in the engine than in
macro code. There is a mix of styles, fonts, nesting, local and per|-|formula
data management involved and not all of that is easy in macros.

Fonts play an important role in the rendering: they provide the symbols and also
the specialized means for composing radicals, fences, scripts, etc. There are two
aspects that have determined how the engine deals with it: the (sometimes pseudo)
width and (sometimes abused) italic correction, extensible snippets that get
combined into a shape of a defined size, and rules that complement these
extensible snippets or stand on their own. None of the fonts that we have
available is perfect and the way around that in \CONTEXT\ is to tweak them at
runtime when we load them.

Tweaks have some history. In \MKIV\ they were mostly responsible for fixing
dimensions and providing additional extensibles. In \LMTX\ we go further and also
use them to improve rendering for which they cooperate with additional engine
features. During the years that we implemented new and better tweaks the
mechanism in \CONTEXT\ that deal with atoms, fractions, radicals, accents and
fences have been partially redone too. Although the engine can do more, we don't
always use the new features. Some are just there because we wanted to improve the
traditional approach. We can actually simplify the math engine a lot by for
instance removing all the code that deals with italic correction (kind of special
in traditional \TEX) and just wipe features we never use.

Among the tasks that the engine delegates to macros is the construction of
adaptive extensibles, especially those combined with text and most of these are
horizontal: arrow or combinations of arrows with other frills, either or not with
tex on top or below. In traditional fonts these are made from tips and rules,
(repeated) minuses, (repeated) equals and other symbols with compensating spaces
in between, to get rid of side bearings, and driven by for instance \TEX's leader
feature. A lot of the chosen methods has to do with the lack of slots in fonts
(to create extensibles) and the fact that commands are used to inject them in
formulas. One can hide a lot in a command (macro) and the user doesn't see how
ugly the hacks are. But no matter what one comes up with, rules are used in
combination with glyphs. Till recently.

When we were making sure that most \OPENTYPE\ math fonts can make nice math using
tweaks, we wondered what to do with the few reasonable \TYPEONE\ math fonts left:
the Antykwa, Iwona and Kurier. In retrospect we should have made a few companion
fonts with extra glyphs but that insight came after re|-|implementing (read:
stripping) virtual \UNICODE\ math to just suit these fonts. A nice challenge was
to make \quotation {everything ruled} to use the more curved rule|-|based symbols
in Antykwa and once that was done a logical next step was to see if that could
also be done with tweaks in \OPENTYPE\ fonts. Once we got started with that the
aim became to completely eradicate the use of rules, which are on the one hand
made for some additions to the engine and on the other hand simplified related
\TEX\ code. It also introduced a few more tweaks for fonts that lack snippets
that replace rules. The mechanisms that are involved are accents (normally
thinner and|/|or smaller symbols), fences (bars), radicals (the top rule) and
what we call stackers in \CONTEXT: the thicker already mentioned arrow like
extensibles. In fact, the smaller accent based arrows were never in \TEX\ fonts
so we cannot use them as stackers: users expect some stability in rendering after
all.

A side effect of this approach is that, if we want, we can now not only strip all
the code that deals with italic correction from the engine, but also all code
related to rules! And we're not only talking of rendering but also all the
variables that somehow determine how rules are done. In the end, when we also rid
code that we don't use half the math engine can go. Of course we've added all
kind of new things too so in the end the size and complexity is still comparable
to the original.

All this doesn't change the fact that the specific math representations that use
some kind of \quote {rulish shape} is still a mess: in \UNICODE\ as well as
\OPENTYPE. First of all the \TEX\ community didn't make sure that the things we
call stackers in \CONTEXT\ got code points, which means that for instance a so
called relbar is basically a stretched minus and that symbols that represent
implying are tagged as left or right arrows. Compare that with the plank constant
h which resulted in omission of the italic math \quote {$h$}.. The same is true
for the vertical bars where there is no distinction between left, right and
middle bars. At the same time the opportunity to get it all right in \OPENTYPE\
fonts has been missed too. That is why we now have to tweak our way out of it.

\starttexdefinition NoMoreRules #1
    \startlinecorrection
    \switchtobodyfont[#1,10pt]%
    \scale[width=\textwidth]\bgroup
        \startframed[offset=1ts,strut=no,width=.5\textwidth]
           \startformula
                   -
             \quad \mrel    {1+2}{a+b+c}
             \quad \overbar {x+y}
             \quad \underbar{x+y}
             \quad \frac    {x+y}{z}
             \quad \sqrt    {x+y}
             \quad -
            \stopformula
        \stopframed
    \egroup
    \blank[nowhite,samepage]
    \midaligned{\strut#1}
    \stoplinecorrection
\stoptexdefinition

\NoMoreRules{antykwa}
\NoMoreRules{bonum}
\NoMoreRules{cambria}
\NoMoreRules{concrete}
\NoMoreRules{dejavu}
\NoMoreRules{ebgaramond}
\NoMoreRules{euleroverpagella}
\NoMoreRules{erewhon}
\NoMoreRules{iwona}
\NoMoreRules{kpfonts}
\NoMoreRules{kurier}
\NoMoreRules{libertinus}
\NoMoreRules{lucida}
\NoMoreRules{modern}
\NoMoreRules{pagella}
\NoMoreRules{schola}
\NoMoreRules{stixtwo}
\NoMoreRules{termes}
\NoMoreRules{xcharter}

From the perspective of the macro package (read: \CONTEXT) it means that we need
to provide extensibles for relbar like stackers, over- and underbars, fraction
separators, and radical top elements. Where possible we can borrow snippets from
the extensible minus and tweak their dimensions and when a font lacks them we
fake extensibles using the regular minus. As a reminder: that ensures that we
have the proper left and right tips, that often have some curvature. In a follow
up we will try to implement a more organic approach to rules (see \in {figure}
[fig:organic-roots]).

\startplacefigure
      [title={An experiment with organic square roots in \TEX.},
       reference=fig:organic-roots]
    \externalfigure[ontarget-rules-001.jpg][width=.8\textwidth]
\stopplacefigure

\stopchapter

\stopcomponent

